\pdfmapfile{+symbols.map}
\pdfmapfile{+cm.map}
\pdfmapfile{+cmextra.map}
\documentclass[11pt]{article}
\usepackage[a4paper,margin=25mm]{geometry}
\usepackage{amsmath,amssymb,amsthm}
\usepackage[colorlinks=true,urlcolor=blue]{hyperref}
\renewcommand{\labelitemi}{$\bullet$}
\newtheorem{theorem}{Theorem}
\newtheorem{lemma}{Lemma}
\newcommand{\E}{\mathbb E}
\newcommand{\C}{\mathcal C}
\title{Equal Density Does Not Give a Constant\\Expected-Subgraph Ratio\\[3pt]\large A counterexample to TLMC 00000006182}
\author{Galaxy-0}
\date{17 October 2026}
\begin{document}
\maketitle
\begin{abstract}
The triangle and the four-cycle are strictly balanced graphs, both of density
one. Their expected numbers of unlabelled, not necessarily induced copies in
$G(n,p)$ have ratio $4/[3(n-3)p]$. Already for $n=4$, the ratio is $16/3$ at
$p=1/4$ and $8/3$ at $p=1/2$. This disproves the conjecture's unqualified
constant-ratio clause. Lean 4.31.0 with only Std checks the actual graph
subsets, strict balance, the complete 64-outcome sample space, symbolic
expectation identities, and a contradiction valid over every
characteristic-zero field. The common-threshold theorem and the narrower
fixed-scale asymptotic statement are not contradicted.
\end{abstract}

\section{Claim, interpretation, and exact scope}
The official English statement ends by asserting that two strictly balanced
subgraphs share a threshold exactly when their densities are equal, ``in
which case the expectation ratio is constant.'' The Chinese statement has
the same unqualified constancy claim. We use the standard model $G(n,p)$:
the vertex set is $\{0,\ldots,n-1\}$, and each unordered pair is independently
an edge with probability $p$. For a fixed graph $H$, let $N_H$ count
\emph{unlabelled subgraphs} isomorphic to $H$. They need not be induced.
Different edge subsets count as different copies, and every copy counts
once, regardless of its automorphisms.

Our counterexample addresses an \emph{exact} ratio claimed to be independent
of the edge probability. In particular it refutes a universal constant
conversion factor obtained solely from equal densities. It does not refute
the ordinary threshold $p\asymp n^{-1}$ for cycles. It also does not refute a
ratio that merely \emph{converges} to a constant after fixing
$p=c/n$ with one fixed $c>0$. Section~\ref{sec:scale} makes the distinction
explicit. If that narrower asymptotic claim was intended, it requires an
additional qualification absent from the source.

\section{Strict balance and the general counting formula}
Write $v(H)$ and $e(H)$ for the numbers of vertices and edges, and
$d(H)=e(H)/v(H)$ for nonempty $H$. A graph is strictly balanced if every
proper subgraph with at least one vertex has smaller density. A subgraph
may discard both vertices and edges; isolated vertices are allowed.

\begin{lemma}
For every $k\geq3$, the cycle $C_k$ is strictly balanced and has density one.
\end{lemma}
\begin{proof}
It has $k$ edges and $k$ vertices. A proper subgraph cannot contain a cycle:
any cycle contained in $C_k$ must traverse all its edges and vertices. Thus
every proper subgraph $J$ is a forest. If $v(J)>0$, then
$e(J)\leq v(J)-1<v(J)$, so $d(J)<1=d(C_k)$.
\end{proof}

\begin{lemma}\label{lem:count}
For $n\geq k\geq3$,
\[
 \E N_{C_k}=\frac{(n)_k}{2k}p^k,
 \qquad (n)_k=n(n-1)\cdots(n-k+1).
\]
\end{lemma}
\begin{proof}
There are $(n)_k$ ordered injective placements of the cyclically ordered
vertices into $[n]$. Each unlabelled cycle has exactly $2k$ descriptions:
a choice of starting vertex and one of the two traversal directions.
Each candidate copy requires $k$ distinct edges, hence is present with
probability $p^k$ by independence. Extra edges have no effect because the
copies need not be induced. Sum the indicator expectations over the
$(n)_k/(2k)$ candidate copies.
\end{proof}

In particular, for all $n\geq4$ and $p>0$,
\begin{align}
 \E N_{C_3}&=\binom n3p^3,&
 \E N_{C_4}&=3\binom n4p^4,\label{eq:general}\\
 \frac{\E N_{C_3}}{\E N_{C_4}}
   &=\frac{4}{3(n-3)p}.\label{eq:ratio}
\end{align}
The factor $3$ in the denominator of~\eqref{eq:ratio} is essential:
the automorphism factors for $C_3$ and $C_4$ are $6$ and $8$.

\section{A finite counterexample with interior probabilities}
\begin{theorem}\label{thm:counter}
The equally dense, strictly balanced graphs $C_3$ and $C_4$ have no constant
expectation ratio in $G(4,p)$ as $p$ varies in $(0,1)$.
\end{theorem}
\begin{proof}
Equation~\eqref{eq:general} gives $\E N_{C_3}=4p^3$ and
$\E N_{C_4}=3p^4$. Exact values are
\[
\begin{array}{c|c|c|c}
 p&\E N_{C_3}&\E N_{C_4}&\E N_{C_3}/\E N_{C_4}\\ \hline
 1/4&1/16&3/256&16/3\\
 1/2&1/2&3/16&8/3
\end{array}
\]
Both denominators are positive, and both probabilities are strictly
between zero and one. Equivalently, the two cross products are
\[
 \frac1{16}\cdot\frac3{16}=\frac3{256}
 \quad\ne\quad
 \frac12\cdot\frac3{256}=\frac3{512}.
\]
Thus no real number $a$ can simultaneously satisfy
$\E N_{C_3}=a\E N_{C_4}$ at both probabilities.
\end{proof}

\section{What happens at the common threshold}\label{sec:scale}
At $p=c/n$, where $c>0$ is fixed, the exact ratio is
\[
 R_n=\frac{4n}{3c(n-3)}
     =\frac4{3c}+\frac4{c(n-3)}.
\]
Consequently $R_n\to4/(3c)$. This correct, narrower asymptotic constancy
is fully compatible with the counterexample. The limiting constant depends
on $c$; even along a fixed scale the finite-$n$ ratio is not exactly constant.
For comparison, $p=n^{-1/2}$ makes the ratio tend to zero, while $p=n^{-2}$
makes it tend to infinity. These conclusions follow immediately from the
exact formula. None requires denying the standard common threshold.

\section{Formal model and audit boundary}
The sole mathematical source file is \texttt{Main.lean}, importing
\texttt{Std}. The formal model is finite but uses actual graphs:
\begin{itemize}
\item Vertices are \texttt{Fin 4}. The six edges are the actual unordered
pairs $01,02,03,12,13,23$, in that order. A graph consists of a four-bit
vertex subset and a six-bit edge subset. Well-formedness requires every
edge endpoint to be present. Subgraphs are inclusion of both subsets.
\item Strict balance is checked against \emph{all} proper subgraphs with
nonempty vertex set, using the equivalent cross-multiplied density
inequality. The chosen graphs have vertex/edge masks $(7,11)$ and $(15,45)$.
\item A three- or four-vertex cycle is defined as a well-formed graph with
the specified number of vertices and degree two at every vertex.
Exhaustive classification gives exactly four triangles and three
four-cycles. A separate theorem proves connectivity using actual walks.
The representative lists have no duplicates, so filtering them by
subgraph inclusion counts copies without automorphism multiplicity.
\item Outcomes are all 64 edge masks, with all four vertices retained.
Every Boolean assignment to the six edges occurs exactly once. The
probability of an outcome with $m$ edges is
$p^m(1-p)^{6-m}$; its equality with the product of its six individual
Bernoulli factors is proved. Total mass is one, and for rational
$0\leq p\leq1$ every mass is nonnegative.
\item Expectation is the finite weighted sum of the actual copy count over
these outcomes. Kernel-checked symbolic identities give $4p^3$ and $3p^4$.
They hold in every field. The final characteristic-zero-field theorem
rules out a common proportionality constant at $p=1/4,1/2$, so the
contradiction applies to real constants, not only rational ones.
\end{itemize}

The complete disproof in Theorem~\ref{thm:counter} is formalized. The
arbitrary-$n$ enumeration proof in Lemma~\ref{lem:count} and the asymptotic
discussion are ordinary mathematical proofs in this report; they are not
claimed as Lean theorems. They provide context and are unnecessary for the
finite counterexample. No external probability library, graph library,
assumed expectation formula, custom axiom, \texttt{sorry},
\texttt{native\_decide}, or kernel-bypassing evaluation is used.

\section{Reproduction and source}
Use Lean 4.31.0, commit
\texttt{68218e876d2a38b1985b8590fff244a83c321783}.
There are no Lake dependencies. From the directory containing this report,
run \texttt{./VERIFY.sh}. It performs a Lake build, re-elaboration with
warnings treated as errors, independent \texttt{leanchecker} kernel replay
of both modules, axiom reports, and a Python standard-library check.
The Python check independently constructs cycles as sets of actual edges,
enumerates all outcomes for $n=4,5$, multiplies Bernoulli factors edge by
edge, and computes exact rational expectations. It is an additional check,
not a source of Lean proof authority. The proof audits report only the
standard logical axioms \texttt{propext}, \texttt{Classical.choice}, and
\texttt{Quot.sound}, as applicable.

The source is
\href{https://github.com/The-Last-Math-Competition/The-Last-Math-Competition/blob/cce2d90cb298de9595798eb674d2c80712100532/conjectures/00000006182.md}{TLMC conjecture 00000006182},
with verified source blob:
\begin{center}\small\texttt{920179a96b1058c56db58613db3647e61f614cd0}.\end{center}
As checked on 17 October 2026, its metadata has neither proven nor
disproven set; the 1,024-item combined historical pull-request snapshot
contains no prior submission referring to this conjecture. Eligibility
and reproduction logs accompany the source package.
\end{document}
